\documentclass[11pt,a4paper]{article}
\usepackage[margin=2.5cm]{geometry}
\usepackage{amsmath,amssymb,booktabs,graphicx,hyperref,natbib}
\usepackage[T1]{fontenc}
\hypersetup{colorlinks=true,linkcolor=blue,citecolor=blue,urlcolor=blue}
\title{\textbf{Energy Shocks and the Green Adoption Channel in a Small Open Economy}}
\author{Guillaume, Boris Chafwehe, Francesca Diluiso}
\date{\today}

\begin{document}
\maketitle

\begin{abstract}
\noindent We add an extensive-margin brown/green durable adoption choice to the
small open economy HANK model of \citet{ARS2024}. Households choose discretely
whether to switch to a green durable that insulates them from fossil energy
prices, paying a resource cost. We study brown energy \emph{price} and
\emph{supply} shocks under the two fiscal responses studied in the recent
literature: a retail price cap \citep{Langot2026} and a Slutsky-compensating
transfer indexed to pre-crisis energy consumption \citep{Bayer2026}. The price
cap stabilises output but mechanically eliminates the green adoption margin,
while the transfer delivers comparable stabilisation and leaves adoption
intact. Shielding consumers from the relative price is therefore not neutral
for the energy transition. The transition loss is linear in the cap's
intensity under a price shock and convex under a supply shock, and the
\emph{sign} of the cap's welfare effect flips with the elasticity of energy
supply --- reconciling \citet{Langot2026} and \citet{Bayer2026} within a
single model. The policy \emph{ordering} is robust to calibration; the
\emph{magnitude} of the adoption channel, and even the sign of its output
effect under a price shock, depend on the calibrated switching cost and taste
scale, which we report and discuss.
\end{abstract}

\section{Introduction}

Three recent papers study fiscal responses to the 2022 European energy crisis
in HANK frameworks. \citet{ARS2024} study a small open economy and stress the
negative externality when all importers subsidise simultaneously.
\citet{Bayer2026} use a two-country HANK with \emph{fixed} euro-area energy
supply and find subsidies are beggar-thy-neighbour while Slutsky transfers are
not. \citet{Langot2026} evaluate the French \emph{bouclier tarifaire} assuming
\emph{perfectly elastic} energy supply and find the price cap performs well.
All three treat the household's energy technology as fixed.

We relax exactly that margin. Households can adopt a green durable that
insulates them from the fossil price, at a resource cost. The energy shock
raises the return to adoption, so the crisis itself can accelerate the
transition --- unless policy removes the price signal that drives it. Our
energy-market closure nests the disagreement between \citet{Bayer2026} and
\citet{Langot2026} in a single parameter: with $\varepsilon^{E}=\infty$ the
economy is a price taker (elastic supply, Langot); with $\varepsilon^{E}$
finite the quantity is fixed and the world price clears the market (Bayer).
Both closures are reported below.

\section{Model}

\subsection{Environment}

The aggregate structure is \citet{ARS2024}: a small open economy with sticky
wages, a CES consumption basket nesting energy against a home/foreign bundle,
imported energy and foreign goods with sticky retail prices, and a mutual fund
holding claims on domestic firms, importers and a share $\zeta^{E}$ of the
energy endowment.

\paragraph{Unit of account.} The household block is numeraire-generic: it
reads a price of the unit of account relative to the CPI, $p^{num}$, a real
return $1+r^{num}=(1+r)\,p^{num}_{-1}/p^{num}$, and after-tax labour income
$atw^{num}=atw/p^{num}$, and returns asset holdings converted by
$A^{CPI}=A\,p^{num}$ before meeting CPI-denominated objects in asset market
clearing and the current account. We set $p^{num}=p_{H}$, the domestic good.
This makes the resource constraint independent of how the CPI weights brown
and green energy, which matters once the two coexist (Section~\ref{sec:cpi}).
Everything contractually written in CPI units --- fiscal transfers, the ARS
subsistence term, and the switching cost $\psi_g$ --- is divided by $p^{num}$
on the way into the budget.

\subsection{The durable state}

Each household carries a durable state $d$, the \emph{pair} (holding entering
the period, holding chosen last period), needed to charge the switching cost:
\[
d\in\{BB,\;BG,\;GB,\;GG\},
\]
where $GB$ (green now, brown before) is the state that pays $\psi_g$. Brown is
absorbing; green breaks down to brown at rate $\delta_g$. The state $BG$
carries zero mass by construction, so the effective state is three-point.
Stages, in forward order, are \texttt{[dep, prod, durables, consav]}:
breakdown, idiosyncratic productivity, the discrete adoption choice (a logit
with taste scale $\sigma_{\varepsilon}$), and the EGM consumption--savings
step. The adoption stage uses the discrete-continuous structure of
\citet{DCEGM2017}.

\subsection{Energy prices}

We distinguish three prices, all relative to the CPI:
\begin{align*}
p^{E} &\quad\text{wholesale/retail market price of brown energy, before any subsidy}\\
P^{E}_{B} &= (1-\tau^{E})\,p^{E}+\tau^{E}\bar p^{E}
&&\text{price paid by a \emph{brown} household}\\
P^{E}_{G} &= \varrho\,\bar p^{E}\left(P^{E}_{B}/\bar P^{E}_{B}\right)^{\rho_g}
&&\text{price paid by a \emph{green} adopter}
\end{align*}
with $\varrho$ the steady-state green/brown ratio and $\rho_g$ the
pass-through of the fossil shock to the green price.

\subsection{How the durable enters the budget}

Households of different durable types face different energy prices, hence
different \emph{consumption bundle prices}. We give each type its own exact
CES cost-of-living index, reusing ARS's own energy-nest parameters
($\alpha_E$, $\eta_E$):
\[
p^{E}_{d}=P^{E}_{B}+\mathbb{I}[d\ \text{green}]\,(P^{E}_{G}-P^{E}_{B}),
\qquad
p_{d}=\Big[1+\alpha_E\big((p^{E}_{d})^{1-\eta_E}-(P^{E}_{B})^{1-\eta_E}\big)\Big]^{\frac{1}{1-\eta_E}},
\]
and the budget constraint is $(p_{d}/p^{num})\,c+a'=\mathrm{coh}(d)$. The
second expression is exact given the CES identity enforced by the equilibrium
block, so brown households have $p_d\equiv 1$ at every date. With
type-specific prices, aggregate energy demand is $\sum_d c_E(d)$, not the CES
demand evaluated at the aggregate index --- the two differ by Jensen's
inequality (0.5\% here) --- so aggregate $c_H$, $c_F$, $c_E$ are built from
the household block's own heterogeneous-agent outputs.

\subsection{What the CPI measures}
\label{sec:cpi}

With two energy prices in the economy, the energy leg of the CPI could be
anchored on the brown price alone or on the share-weighted index
$[(1-D^{G})(P^{E}_{B})^{1-\eta_E}+D^{G}(P^{E}_{G})^{1-\eta_E}]$. We anchor on
the brown price: it is what statistical agencies do (the energy basket is not
reweighted quarterly for technology adoption), it preserves $p_{d}\equiv 1$
for brown households exactly, and feeding $D^{G}$ back into the CPI would
close a DAG cycle requiring the CPI to be promoted to a model unknown. Because
the unit of account is the domestic good, this choice does not touch the
resource constraint; it affects only the measured deflator and, through it,
the monetary rule. The gap between the share-weighted and brown-anchored CPI
has closed form
\[
\phi^{1-\eta_E}=1+\alpha_E\,D^{G}\Big[(P^{E}_{G})^{1-\eta_E}-(P^{E}_{B})^{1-\eta_E}\Big],
\]
computed ex post outside the model. On the price shock without policy it
reaches 0.14 annualised percentage points, about 1.5\% of the amplitude of
measured CPI inflation (0.23 pp, 2.3\%, on the supply shock; Figure~\ref{fig:deflator}).

\begin{figure}[h]\centering
\includegraphics[width=0.7\textwidth]{output/fig8_deflator_gap.png}
\caption{Measurement gap from anchoring the CPI on $P^{E}_{B}$ alone: true
minus measured annualised inflation, computed ex post.}
\label{fig:deflator}
\end{figure}

\subsection{Balance of payments}
\label{sec:bop}

Two durable-margin terms enter the current account. The \textbf{switching
cost} is booked as an import, $\text{imports}_{dur}=\psi_g D^{sw}$: booking it
as domestic demand instead fails here, since with $\mu_{ss}>1$ and
capitalised profits, extra demand for the home good creates a marginal
dividend owned by nobody. The \textbf{energy saving} of adopters,
$(P^{E}_{B}-P^{E}_{G})\,C^{G}_{E}$, is booked as a real import saving: this
asserts that cheaper energy for adopters is a genuine resource saving for the
country, not a domestic transfer, and gives the windfall an external
counterpart consistent with asset market clearing.

\subsection{Policy instruments}

Both instruments are the extreme (full-compensation) versions, as in
\citet{Bayer2026}.
\begin{align}
\text{Price cap:}\quad & P^{E}_{B,t}=(1-\tau^{E})p^{E}_{t}+\tau^{E}\bar p^{E},
&& \tau^{E}=1 \Rightarrow \text{retail price fixed at pre-crisis level}\\
\text{Slutsky transfer:}\quad & tr_{it}=\iota\,(p^{E}_{t}-\bar p^{E})\,\bar c_{E,i},
&& \iota=1 \Rightarrow \text{full compensation on \emph{pre-crisis} use}
\end{align}
Both are debt financed and repaid through the distortionary labour tax. The
transfer base $\bar c_{E,i}$ is each household's steady-state energy demand.
The cap operates directly on $P^{E}_{B}$, exactly the price entering the
brown--green gap, and therefore compresses the adoption incentive; the
transfer is lump-sum with respect to current choices and leaves the relative
price intact.

\section{Calibration}

\begin{table}[h]\centering
\small
\begin{tabular}{lll}
\toprule
Parameter & Value & Source / target \\
\midrule
$\alpha_E$ energy expenditure share & 0.04 & \citet{ARS2024} \\
$\eta_E$ energy substitution & 0.10 & \citet{ARS2024} \\
$r$, $\mu_{ss}$ & 0.01, 1.03 & \citet{ARS2024} \\
$\beta$ heterogeneity (3 types) & spread 0.06 & MPC $\approx0.30$ \\
\midrule
$\delta_g$ green breakdown rate & 0.05 & assumption (5\,yr life) \\
$\varrho=P^{E}_{G}/P^{E}_{B}$ steady-state gap & 0.80 & assumption \\
$\psi_g$ switching cost & 0.538 & targets $D^{G}_{ss}=0.05$ \\
$\sigma_{\varepsilon}$ taste scale & 0.05 & calibrated jointly with $\psi_g$ \\
$\rho_g$ pass-through to green price & 0 & upper bound on the channel \\
\bottomrule
\end{tabular}
\caption{Calibration. The steady state is identical across all experiments.
$\psi_g$ and $\sigma_\varepsilon$ are calibrated jointly to the level target
$D^{G}_{ss}=0.05$; $\sigma_\varepsilon$ is not separately pinned down by an
adoption-elasticity moment (Section~\ref{sec:limits}).}
\end{table}

\section{Experiments and results}

The \emph{price} shock is a 100 log-point rise in the world energy price with
a 16-quarter half-life, as in \citet{ARS2024}. The \emph{supply} shock is a 5\%
cut in available energy for 6 quarters, after which the world price clears
the market endogenously. Impulse responses are reported as $100\times$ the
level deviation from steady state: a percent deviation for variables
normalised to one ($y$, $C$, $n$, $w$, prices), percentage points for
population shares ($D^{G}$, $D^{sw}$), and a share of the energy steady state
($\approx0.040$) for energy quantities.

\begin{table}[h]\centering
\small
\begin{tabular}{llrrrrrr}
\toprule
Shock & Policy & $y(0)$\% & $\sum y$ & $\pi(0)$ ann. & peak $D^{G}$ pp & fiscal cost & adoption $\to \sum y$\\
\midrule
price  & none     & $-1.12$ & $-30.1$ & 9.14  & \textbf{9.18} & 0.00 & $-1.60$\\
price  & cap      & $-0.00$ & $-20.9$ & $-1.53$ & \textbf{0.01} & 53.82 & $+0.36$\\
price  & transfer & $+1.20$ & $+12.7$ & 11.08 & \textbf{9.32} & 53.39 & $-1.74$\\
\midrule
supply & none     & $-0.71$ & $-8.16$ & 10.06  & \textbf{2.33} & 0.00 & $+0.38$\\
supply & cap      & $-0.74$ & $-46.9$ & $-0.96$ & \textbf{0.01} & 26.09 & $+0.18$\\
supply & transfer & $+1.33$ & $+2.72$ & 14.46 & \textbf{3.39} & 26.56 & $+2.48$\\
\bottomrule
\end{tabular}
\caption{Cumulative sums over 24 quarters. The last column is the contribution
of the adoption margin: the difference between the model with the margin open
and the same model with it shut.}
\label{tab:summary}
\end{table}

\subsection{The crisis moves the transition, in a shock-dependent direction}

With no policy, the green share rises by up to 9.2 percentage points under
the price shock (from 5.0\% to 14.2\% of households) and 2.3 pp under the
supply shock (Figure~\ref{fig:baseline}). The output effect of the adoption
margin, however, is not one-signed. Under the \textbf{supply} shock it is
expansionary ($+0.38$ points of cumulative output with no policy, $+2.48$
under the transfer): adopters escape the energy bill and their real income
falls less. Under the \textbf{price} shock it is contractionary ($-1.60$
points with no policy, $-1.74$ under the transfer): the switching cost is
booked as an import in exactly the periods household income is falling
hardest (Section~\ref{sec:bop}), and at the calibrated $\psi_g$ this resource
cost outweighs the energy-bill saving. The sign of the adoption channel's
contribution to output is therefore not a general property of the model --- it
depends on the calibrated switching cost and on which shock is driving
adoption.

\begin{figure}[h]\centering
\includegraphics[width=0.72\textwidth]{output/fig1_price_shock.png}
\includegraphics[width=0.72\textwidth]{output/fig1_supply_shock.png}
\caption{Brown energy shocks, adoption open (solid) vs.\ shut (dashed), no
policy. Top: price shock, elastic energy supply. Bottom: supply shock, $-5\%$
for 6 quarters, fixed quantity.}
\label{fig:baseline}
\end{figure}

\subsection{The price cap eliminates the adoption margin}

Under the cap, the peak green share response falls to essentially zero under
both shocks (from $+9.18$ to $+0.01$ pp on the price shock, from $+2.33$ to
$+0.01$ on the supply shock): the adoption channel is not attenuated, it is
switched off. Its contribution to cumulative output collapses to near zero
accordingly ($-1.60\to+0.36$ on the price shock, $+0.38\to+0.18$ on the supply
shock). The mechanism is direct: the cap holds $P^{E}_{B}$ at its pre-crisis
level, and $P^{E}_{B}$ is precisely what determines the return to adoption.

\begin{figure}[h]\centering
\includegraphics[width=0.72\textwidth]{output/fig2_policy_price.png}
\includegraphics[width=0.72\textwidth]{output/fig2_policy_supply.png}
\caption{Fiscal response, adoption margin open: no policy (solid), price cap
(dashed), Slutsky transfer (dotted). Top: price shock. Bottom: supply shock.}
\end{figure}

\subsection{The transfer preserves adoption, at a cost to output}

The Slutsky transfer leaves the green share response essentially unchanged
($9.32$ vs.\ $9.18$ pp under the price shock, $3.39$ vs.\ $2.33$ under the
supply shock) while delivering more output stabilisation than the cap, at
similar fiscal cost: subsidies destroy the transition margin, transfers do
not. Leaving adoption intact is not free for output under the price shock,
however: the transfer's cumulative-output gain is \emph{smaller} with
adoption open than shut ($-1.74$ points from the margin, Table~\ref{tab:summary}),
for the same resource-cost reason as above.

\subsection{Fixed supply makes the cap destructive}

Under the supply shock the cap yields cumulative output of $-46.9$ against
$-8.2$ with no policy: with inelastic supply, capping the price prevents
substitution away from energy and pushes the economy towards Leontief. This
survives the addition of an adoption margin. It is also precisely the
configuration in which \citet{Langot2026}, assuming elastic supply, find the
cap effective: our elastic-supply column agrees with them ($-20.9$ vs.\
$-30.1$ for cap vs.\ no policy under the price shock).

\begin{figure}[h]\centering
\includegraphics[width=\textwidth]{output/fig3_policy_adoption.png}
\caption{Policy and the adoption margin. Under both shocks the price cap
(dashed) flattens the green share response to zero; the transfer (dotted)
leaves it intact.}
\end{figure}

\subsection{The cap is a dial, not a switch}

Table~\ref{tab:summary} used $\tau^{E}=1$, a complete price freeze. Table~\ref{tab:dose}
sweeps $\tau^{E}\in[0,1]$. Under the \emph{price} shock the adoption loss is
almost exactly linear in $\tau^{E}$ (half a cap costs a bit under half the
transition: 4.53 against 9.18 pp). Under the \emph{supply} shock it is convex:
small caps barely touch adoption and the collapse is concentrated near the
full freeze. A partial shield is therefore much less damaging to the
transition when supply is inelastic --- but, as the next result shows, much
more damaging to everything else.

\begin{table}[h]\centering\small
\begin{tabular}{llrrrrr}
\toprule
Shock & Policy & peak $D^{G}$ pp & $\sum y$ & $\pi(0)$ & fiscal & CEV \%\\
\midrule
price & cap $\tau^E=0$    & 9.18 & $-30.1$ & 9.14 & 0.00 & $-1.30$\\
price & cap $\tau^E=0.25$ & 6.86 & $-27.8$ & 6.50 & 13.39 & $-1.12$\\
price & cap $\tau^E=0.50$ & 4.53 & $-25.5$ & 3.84 & 26.83 & $-0.94$\\
price & cap $\tau^E=0.75$ & 2.20 & $-23.2$ & 1.17 & 40.30 & $-0.76$\\
price & cap $\tau^E=1$    & 0.01 & $-20.9$ & $-1.53$ & 53.82 & $-0.57$\\
price & transfer          & 9.32 & $+12.7$ & 11.08 & 53.39 & $-0.21$\\
\midrule
supply & cap $\tau^E=0$    & 2.33 & $-8.16$ & 10.06 & 0.00 & $-0.59$\\
supply & cap $\tau^E=0.25$ & 2.21 & $-10.60$ & 9.28 & 5.96 & $-0.63$\\
supply & cap $\tau^E=0.50$ & 1.99 & $-14.80$ & 8.07 & 15.34 & $-0.70$\\
supply & cap $\tau^E=0.75$ & 1.50 & $-23.27$ & 5.89 & 30.89 & $-0.84$\\
supply & cap $\tau^E=1$    & 0.01 & $-46.85$ & $-0.96$ & 26.09 & $-1.28$\\
supply & transfer          & 3.39 & $+2.72$  & 14.46 & 26.56 & $-0.14$\\
\bottomrule
\end{tabular}
\caption{Dose-response in the cap intensity. CEV is the consumption-equivalent
variation relative to the steady state; negative means households would pay
that share of permanent consumption to avoid the scenario.}
\label{tab:dose}
\end{table}

\subsection{The welfare effect of the cap flips sign with the supply elasticity}

This is the sharpest result in the paper. Under the price shock, tightening
the cap \emph{raises} welfare monotonically, from $-1.30\%$ to $-0.57\%$.
Under the supply shock it \emph{lowers} welfare monotonically, from $-0.59\%$
to $-1.28\%$. \citet{Langot2026} assume elastic supply and conclude the French
shield worked; \citet{Bayer2026} assume fixed supply and conclude subsidies
are destructive. Both are right within their own closure, and our model
reproduces each by moving a single parameter: the disagreement in the
literature is not about mechanisms, it is an empirical question about the
elasticity of energy supply at the relevant horizon.

\begin{figure}[h]\centering
\includegraphics[width=\textwidth]{output/fig6_dose_response.png}
\caption{Dose-response in $\tau^{E}$ (orange) against the Slutsky transfer
benchmark (dotted). Note the sign reversal of the welfare panel between the
two shocks.}
\end{figure}

\subsection{Distributional incidence}

\begin{table}[h]\centering\small
\begin{tabular}{lrrrr}
\toprule
Scenario (price shock) & CEV mean & impatient & middle & patient\\
\midrule
no policy        & $-1.30$ & $-2.10$ & $-1.65$ & $-0.14$\\
cap $\tau^E=0.5$ & $-0.94$ & $-1.45$ & $-1.25$ & $-0.11$\\
cap $\tau^E=1$   & $-0.57$ & $-0.78$ & $-0.85$ & $-0.08$\\
transfer         & $-0.21$ & $-0.36$ & $-0.45$ & $+0.17$\\
\bottomrule
\end{tabular}
\caption{CEV by discount-factor type (\%). Impatient types hold little wealth
and have high MPCs.}
\end{table}

Without policy the impatient lose far more than the patient ($-2.10$ vs.\
$-0.14$, roughly $15\times$). The patient remain the least-hurt group
throughout the cap sweep. One clear distributional result: under the
transfer, the \textbf{patient type's CEV turns positive} ($+0.17\%$) while
impatient and middle types remain negative. Indexed to pre-crisis energy use,
the Slutsky compensation over-compensates low-MPC, low-energy-exposure
households relative to their true loss --- a symptom of the fiscal block's
lack of investment margin to crowd out the transfer (Section~\ref{sec:limits}).

\section{Solution accuracy}
\label{sec:accuracy}

Solving the model under the CPI numeraire instead of the domestic good gives a
bit-identical steady state and impulse responses agreeing to $2\times10^{-5}$
on $y$ and $3\times10^{-5}$ on $C$ (about $0.15\%$ of the respective
amplitudes) --- a relabeling of the same economy, not two different models.

The first-order solution understates the adoption margin at the shock sizes
used in this literature. Table~\ref{tab:nonlinear} compares the fully
nonlinear general-equilibrium transition against its linearised counterpart
at several price-shock sizes.

\begin{table}[h]\centering\small
\begin{tabular}{lrrrrr}
\toprule
Non-linear / linear & $y$ & $C$ & $c_E$ & $D^{G}$ & $\pi$ ann.\\
\midrule
size 0.125 & 0.96 & 0.97 & 0.94 & 1.16 & 0.97\\
size 0.25  & 0.91 & 0.93 & 0.87 & 1.33 & 0.95\\
size 0.50  & 0.78 & 0.82 & 0.73 & 1.67 & 0.92\\
\bottomrule
\end{tabular}
\caption{Ratios at the date of the linear peak. Output and inflation are
damped relative to the linear prediction while $D^{G}$ is amplified: a larger
share of households escapes the energy bill than the linear solution
predicts. The nonlinear solver does not converge at size 1.00, the size used
in Table~\ref{tab:summary} for the price shock.}
\label{tab:nonlinear}
\end{table}

The 5\% supply cut used throughout this paper (rather than the 20\% cut used
in some of the literature) is motivated by the same non-convergence concern:
the first-order solution is not reliable at larger shock sizes, and at some
sizes there is no nonlinear benchmark to compare it against at all.

\section{Limitations}
\label{sec:limits}

\begin{enumerate}\itemsep3pt

\item \textbf{The taste scale $\sigma_{\varepsilon}$ is calibrated, not
estimated.} $\psi_g$ and $\sigma_\varepsilon$ jointly hit $D^G_{ss}=0.05$, but
one level moment cannot separately identify both a switching cost and a taste
scale: a second moment (e.g.\ the semi-elasticity of adoption to a purchase
subsidy) is needed and would discipline the sign found in
Section~4.1 along with the magnitudes throughout. The \emph{ordering} of the
cap-vs-transfer results does not depend on it, since it follows from the
relative price alone.

\item \textbf{The first-order solution understates the adoption margin} at
the shock sizes used in this literature (Section~\ref{sec:accuracy}), which
in turn overstates the recession it causes.

\item \textbf{Zero pass-through to the green price} ($\rho_g=0$) means
adopters are completely insulated from the fossil shock: an upper bound on
the adoption channel. A partial pass-through calibrated to observed
electricity--gas co-movement would be more defensible.

\item \textbf{No within-country energy-share heterogeneity.}
\citet{Bayer2026}'s central distributional mechanism is that poor households
have higher energy shares; we have homothetic energy demand conditional on
the durable type, so we cannot speak to their distributional results.

\item \textbf{No investment margin.} Both \citet{Bayer2026} and
\citet{ARS2024} have richer supply sides. Its absence here likely makes the
transfer look more expansionary, and the supply shock more inflationary, than
it would be with capital to smooth the adjustment.

\end{enumerate}

\begin{thebibliography}{9}

\bibitem[Auclert et al.(2024)]{ARS2024}
Auclert, A., H. Monnery, M. Rognlie, and L. Straub (2024).
``Managing an Energy Shock: Fiscal and Monetary Policy.''
In \emph{Heterogeneity in Macroeconomics: Implications for Monetary Policy},
Series on Central Banking, Analysis, and Economic Policies, Vol. 30,
Central Bank of Chile, 39--108.

\bibitem[Bayer et al.(2026)]{Bayer2026}
Bayer, C., A. Kriwoluzky, G. J. M\"uller, and F. Seyrich (2026).
``Redistribution Within and Across Borders: The Fiscal Response to an Energy
Shock.'' \emph{Journal of Political Economy Macroeconomics}, forthcoming.
Earlier version: ``Hicks in HANK,'' CEPR DP 18557.

\bibitem[Iskhakov et al.(2017)]{DCEGM2017}
Iskhakov, F., T. H. J{\o}rgensen, J. Rust, and B. Schjerning (2017).
``The Endogenous Grid Method for Discrete-Continuous Dynamic Choice Models with
(or without) Taste Shocks.'' \emph{Quantitative Economics} 8(2), 317--365.

\bibitem[Langot et al.(2026)]{Langot2026}
Langot, F., S. Malmberg, F. Tripier, and J.-O. Hairault (2026).
``The Macroeconomic and Redistributive Effects of the French Tariff Shield.''
\emph{Journal of Political Economy Macroeconomics}, forthcoming.
Earlier version: CEPREMAP Working Paper 2305.

\end{thebibliography}

\end{document}
